% Asymmetric Judge Error and Clustered Frames.
%
% Numeric tables are NOT written here. They are generated from the committed result files by
% scripts/export_paper_tables.py into paper/generated/ and pulled in with \input, so a table
% cannot drift from the artifact it reports. Figures quoted in prose are pinned by
% scripts/check_prose_figures.py, the same way the two markdown writeups are (D081, D086).
\documentclass[11pt]{article}

\usepackage[margin=1.1in]{geometry}
\usepackage{booktabs}
\usepackage{amsmath}
\usepackage{microtype}
\usepackage[hidelinks]{hyperref}
\usepackage[numbers,sort&compress]{natbib}

\setlength{\tabcolsep}{4pt}
\sloppy

\newcommand{\vernierGitRev}{b7f2867}
\newcommand{\vernierCardDigest}{f850d3622c13d64c}
\newcommand{\vernierNPrimary}{153}
\newcommand{\vernierNClaims}{18}
\newcommand{\vernierVerdict}{NOT\_VERIFIED}


\title{Asymmetric Judge Error and Clustered Frames:\\
A Pre-Registered Audit of a Commercial Dataset Quality Claim}
\author{Caio Theodoro\\\small Independent}
\date{\today}

\begin{document}
\maketitle

\begin{abstract}
\noindent
Commercial dataset vendors increasingly advertise corpus quality with statistics produced by a
vision-language model applied once to a sample of frames. We audit one such claim. Build AI
publishes, for its Egocentric-10K corpus, that 96.42\% of frames show at least one hand and
91.66\% show active manipulation, both ahead of Ego4D and EPIC-KITCHENS-100, from 10{,}000
frames labelled once by a single model with no human label, no interval, and no test that the
labeller scores the compared corpora on the same scale. We pre-registered a protocol, replicated
the figures with an independent open-weights judge on two releases a year apart, collected
\vernierNPrimary{} human labels against a written rubric, and re-estimated the published
quantities with prediction-powered inference. Three findings. First, judge error against human
gold is strongly asymmetric: of 28 disagreements, 26 inflate the statistic, and on hand counting
there are 16 over-reports against zero under-reports across 129 frames that could have been
under-reported. The vendor's own stored labels, from a different model family, carry the same
sign. Second, the corpus is clustered by camera wearer, with a design effect of 1.25 to 1.66
measured on two independent 10{,}000-frame draws, so an interval computed as if frames were
independent is 12 to 29\% too narrow; no interval is published, and the released evaluation
frames carry no identifier that would let anyone else compute one. Third, the comparative claims
cannot be resolved at this gold size in either direction: correcting both sides of the
manipulation margin against EPIC-KITCHENS-100 moves it from a published $+6.62$ to $-1.02$
percentage points with a 95\% interval of $[-11.68, +9.65]$, which covers the published value.
We report what it would cost to settle it, and we report that buying more labels moved the
estimate toward the vendor and raised that cost. Artifacts, protocol and a decision log are
public; the measurement card reports \texttt{\vernierVerdict} because one pre-registered item
could not be checked at all.
\end{abstract}

\section{Introduction}

A dataset vendor sells a corpus and a number attached to it. Build AI sells egocentric video
recorded in factories, and its evaluation release publishes the quality figures that justify the
sale: 96.42\% of frames contain at least one hand, 76.34\% contain two, 91.66\% show active
manipulation, each ahead of the two most-used academic egocentric corpora. The corpus is
substantial and the release is unusually open by the standards of commercial data: 10{,}000 hours
from 2{,}153 workers across 85 factories, per-clip metadata, camera intrinsics, and the
evaluation frames themselves.

The number attached to it is a point estimate produced by applying \texttt{gemini-2.5-flash} once
to 10{,}000 frames. It carries no human label, no confidence interval, and no evidence that the
model scores a factory floor and a domestic kitchen on the same scale. The last of these is what
the comparative claim rests on, and the margin over EPIC-KITCHENS-100 is about six percentage
points on both headline tasks. A buyer has no way to check any of it.

We built the check. This paper reports what it found and, equally, what it could not settle.

\paragraph{Contributions.}
\begin{enumerate}
\item \textbf{An asymmetric-error result.} On \vernierNPrimary{} human-labelled frames the
      independent judge makes 28 errors and 26 of them inflate the published quantity. On hand
      counting the asymmetry is absolute at this sample size: 16 over-reports, zero
      under-reports. The vendor's own stored labels, produced by a different model family,
      carry the same sign. A judge-derived prevalence of this kind is therefore better read as
      an upper bound than as a point estimate.
\item \textbf{A design effect for a corpus quality statistic.} Frames from one camera wearer
      share scene, lighting, task and equipment. Measured on two independent 10{,}000-frame
      draws from the raw corpus with real worker identifiers, the design effect is 1.25 to 1.66,
      so an interval that treats frames as independent is 12 to 29\% too narrow. We are not
      aware of an egocentric dataset that reports this, and the released evaluation frames carry
      a bare UUID with no worker linkage, so it cannot be computed on the sample the published
      figures describe.
\item \textbf{A negative result on the comparative claim, with a price.} Correcting both sides
      of the margin leaves the published value inside the corrected interval. We report the
      gold-set size that would separate them, and that this figure rose after we collected more
      labels.
\item \textbf{An audit that reports its own failures.} Five corrections to our own published
      claims are recorded in Section~\ref{sec:moved}, including a protocol violation we found
      only by measuring it, a rubric rule nothing enforced, and a coverage table that claimed an
      experiment we never ran.
\end{enumerate}

\section{Related work}

\paragraph{Gold-corrected estimation with model labels.}
Prediction-powered inference \citep{angelopoulos2023ppi} and its efficient variant
\citep{angelopoulos2023ppiplusplus} produce valid intervals from a small human-labelled sample
plus a large model-labelled one, with no assumption that the model is unbiased. Confidence-driven
inference \citep{gligoric2025cdi} refines the same idea by spending human effort where the model
is least confident. A 2026 line applies this machinery to language-model evaluation:
\citet{divekar2026ranking} extends PPI to ranking metrics such as Precision@$K$, and
\citet{feng2026noisybutvalid} estimates a judge's true- and false-positive rates from a small
calibration set in order to certify that a model's failure rate sits below a threshold.

Our estimand differs from all of these in a way that matters. Each of them estimates a property
of \emph{a model, a ranking, or a system}, on an evaluation the authors themselves constructed.
We estimate a property of \emph{a corpus}, published by \emph{a third party}, using that party's
own released frames and their own stored labels, and we ask whether their published number
survives the correction. We are not aware of prior work that does this.

The connection to \citet{feng2026noisybutvalid} is worth stating precisely, because it frames our
first result. Their method requires both error rates, since a judge that misses and a judge that
over-calls distort an estimate differently. In our data one of the two rates is at or near zero.
That is a special case of their setting rather than a challenge to it, and it is the case in
which a naive judge-only estimate is biased in a knowable direction.

\paragraph{Judge validation and its limits.}
Agreement between a model judge and a human is routine to report and easy to report badly.
\citet{rao2026agreement} shows that protocol choices alone move a reported agreement figure
substantially and proposes a reporting checklist; we follow it, and we report Gwet's AC1
\citep{gwet2008ac1} alongside Cohen's $\kappa$ because $\kappa$ is uninformative at the base
rates involved here, where one class exceeds 90\%.

\paragraph{Clustered inference in evaluation.}
\citet{miller2024errorbars} argues that evaluation results deserve error bars and that
dependence between evaluation items inflates the standard errors that flat-rate reporting
ignores. Our design-effect measurement is that argument applied to a dataset statistic rather
than a model score, with the clustering variable being the person wearing the camera.

\paragraph{Prompt sensitivity.}
\citet{liu2026ipr} formalises inter-prompt reliability and its pairwise agreement rate for
model-based labelling. Our prompt sweep reports the same quantities.

\paragraph{Selective prediction.}
\citet{jung2024trustorescalate} guarantees a user-specified level of human agreement by
abstaining, with thresholds selected under distribution-free risk control
\citep{angelopoulos2021ltt}. Our distillation arm targets that guarantee and does not reach it.

\paragraph{Agreement practice in egocentric datasets.}
HD-EPIC \citep{perrett2025hdepic} is the most rigorous annotation-agreement pipeline we found in
this area, with multiple annotators per narration and a minimum agreement threshold, but it
governs the timing of action segments rather than the density of a model-produced attribute, and
it reports no interval on a dataset-level number. EPIC-KITCHENS-100 \citep{damen2022epic100}
reports annotation-density improvements and no inter-annotator agreement statistic. Across the
egocentric datasets whose full text we read, none reports a confidence interval or an agreement
statistic on a headline quality claim. We did not exhaust the field, and we record that as a
best-effort absence rather than a proof.

\section{The audited claim, and what the release ships}

The evaluation release contains three parquet files, one per corpus, 10{,}000 frames each, with
the image bytes and the vendor's own per-frame labels. That is a genuinely useful release: the
comparison can be run on the vendor's own frames, so no sampling difference can explain a
disagreement.

Three properties of it bound what any audit can do.

\paragraph{The frames cannot be traced to the corpus.} \texttt{frame\_id} is a bare UUID4 with no
factory, worker, clip or timestamp component. The published sample therefore cannot be checked
against the corpus it was drawn from by anyone outside the vendor, and no clustered interval can
be computed on it. Our design effect is measured on our own draws for this reason.

\paragraph{The identifier was present and is no longer.} The 10K evaluation release and the 100K
release that supersedes it differ in schema: the later one does not carry the \texttt{frame\_id}
column the earlier one shipped. We state this as a fact about reproducibility and draw no
inference about why.

\paragraph{The corpus count differs slightly from the published one.} Indexing all 19{,}495
shards of the raw corpus yields 2{,}144 distinct workers against a published 2{,}153, with
factories, hours and clip count reconciling almost exactly. Nine workers, 0.42\%, consistent with
ordinary release-time filtering and too few to affect any figure here.

\section{Protocol}

The protocol was frozen before any frame was fetched and every departure from it is a dated entry
in a public decision log, currently 85 entries.

\paragraph{Judge.} \texttt{Qwen/Qwen3-VL-8B-Instruct-FP8}, self-hosted, temperature 0 and seed
777 pinned on every call, with the vendor's prompts used verbatim. The original judge could not
be used: \texttt{gemini-2.5-flash} is deprecated for new API keys, which we confirmed with a live
404. That is itself a reproducibility finding, since a quality guarantee whose instrument is a
versionless third-party API cannot be re-run by its own vendor. The comparison therefore runs
against the vendor's stored per-frame labels, which ship in the release.

\paragraph{Human gold.} One rater labelled \vernierNPrimary{} frames against a written rubric of
twelve numbered rules resolving what the vendor's prompt leaves open, such as gloves,
reflections, motion blur, and an idle grip that is not manipulation. The labelling tool has no
code path to judge output, so the rater cannot see the model's answer.

\paragraph{Estimators.} Prevalence by PPI++ \citep{angelopoulos2023ppiplusplus}; agreement by
Gwet's AC1 \citep{gwet2008ac1} with $\kappa$ beside it; intervals by cluster bootstrap over
worker identifier, $B = 10{,}000$, wherever such an identifier exists, and by an iid bootstrap
labelled as a lower bound on width where it does not.

\paragraph{The full slate, before any result.} Table~\ref{tab:ledger} lists every pre-registered
hypothesis with its threshold and outcome. We put it here, ahead of the results, so that the
reader sees the failures at the same time as the findings.

\begin{table}[t]
\centering\small
\begin{tabular}{@{}lp{0.34\linewidth}p{0.15\linewidth}p{0.20\linewidth}l@{}}
\toprule
id & pre-registered claim & threshold & observed & verdict \\
\midrule
H8 & participant counts differ across corpora & computable, no threshold & 58.2$\times$ & reported \\
H1 & all three figures reproduce within $\pm$2\,pp & $\pm$2\,pp & 6.32\,pp on 2 hands & fails \\
H1b & the two P0 prompts disagree by $\geq$1\,pp & $\geq$1\,pp & 0.32\,pp & null \\
H2 & design effect over worker id $\geq$ 2 & $\geq$ 2 & 1.66 max & fails \\
H3 & manipulation prompt spread $\geq$ 5\,pp & $\geq$ 5\,pp & 1.25\,pp & fails \\
H4 & AC1 higher on hand count than manipulation & direction & 0.860 vs 0.887 & reversed \\
H5 & judge error higher on EPIC by $\geq$ 5\,pp & $\geq$ 5\,pp & +3.57\,pp & fails \\
H6 & agreement floor 0.80 at coverage $\geq$ 0.70 & 0.80 at 0.70 & fidelity 0.693 & fails \\
H7 & calibration not measurable under the published protocol & n/a & ECE 0.104 & measured, degenerate \\
\midrule
\multicolumn{5}{l}{\emph{Not pre-registered, reported as exploratory:}} \\
--- & gold-corrected margin vs.\ EPIC-KITCHENS-100 & none & -1.02\,pp & unresolved \\
\bottomrule
\end{tabular}

\caption{Every pre-registered hypothesis, its threshold, and its outcome. Four of the eight fail
against their own criteria and one reverses direction. The exploratory row was not
pre-registered and is reported as such.}
\label{tab:ledger}
\end{table}

\section{Replication}

Table~\ref{tab:replication} compares the independent judge against the published figures on both
releases at $N = 10{,}000$ per release.

\begin{table}[t]
\centering\small
\begin{tabular}{llrrr}
\toprule
release & figure & measured & published & diff (pp) \\
\midrule
Egocentric-10K & $\geq$1 hand & 95.45 & 96.42 & 0.97 \\
Egocentric-10K & 2 hands & 82.66 & 76.34 & 6.32$^{\dagger}$ \\
Egocentric-10K & active manipulation & 91.28 & 91.66 & 0.38 \\
Egocentric-100K & $\geq$1 hand & 96.09 & 96.95 & 0.86 \\
Egocentric-100K & 2 hands & 85.19 & 79.05 & 6.14$^{\dagger}$ \\
Egocentric-100K & active manipulation & 92.14 & 92.76 & 0.62 \\
\bottomrule
\end{tabular}

\caption{Independent replication on both releases. $\dagger$ marks a figure outside the
pre-registered $\pm$2\,pp tolerance.}
\label{tab:replication}
\end{table}

Hand visibility and active manipulation land within about a point of the published values on both
releases, a year apart. The two-hands figure does not, and misses in the same direction both
times.

We do not know why, and the parsimonious explanation implicates our judge rather than the
vendor's number: on the gold set our judge calls two hands more often than the rater does. What
keeps the observation alive is that the same figure carries the largest design effect on both
corpus draws (Section~\ref{sec:cluster}). Three measurements single out one column and none of
them explains it. We report it and stop there.

\section{Judge error against human gold}

\begin{table}[t]
\centering\small
\begin{tabular}{lrrrr}
\toprule
task & over & at risk & under & at risk \\
\midrule
hand count & 16 & 59 & 0 & 129 \\
active manipulation & 10 & 33 & 2 & 120 \\
\bottomrule
\end{tabular}

\caption{Direction of judge error against the rater. ``At risk'' counts the frames on which the
error was possible: a frame the rater scored as two hands cannot be over-counted, and one scored
as zero cannot be under-counted.}
\label{tab:errdir}
\end{table}

Table~\ref{tab:errdir} is the paper's first result. Of 28 disagreements, 26 inflate the
statistic. On hand counting the asymmetry is absolute at this sample size, 16 against zero, and
the zero is across 129 frames where an under-report was possible. On manipulation it is 10
against 2.

Two checks on whether this is an artifact of our rubric or our judge.

The rubric's most discretionary rule is the one excluding an idle grip from manipulation, and it
would be the natural source of manufactured disagreement. None of the manipulation false
positives carries the corresponding tag. They are tagged for blur, self-contact, another person
in frame, or nothing at all.

The vendor's own stored labels, from a different model family, disagree with the rater in the
same direction on the same frames. Whatever produces the asymmetry is not specific to the model
we chose.

\section{Corrected prevalence, and the comparative claim}

\begin{table}[t]
\centering\small
\begin{tabular}{@{}llrrrlr@{}}
\toprule
corpus & task & published & judge & PPI++ & 95\% CI & $n_{\text{gold}}$ \\
\midrule
Egocentric-10K & $\geq$1 hand & 96.42 & 95.50 & 92.43 & [86.95, 97.92] & 63 \\
Egocentric-10K & active manipulation & 91.66 & 90.00 & 85.06 & [77.90, 92.21] & 63 \\
Ego4D & $\geq$1 hand & 67.33 & 68.00 & 61.20 & [50.22, 72.19] & 30 \\
Ego4D & active manipulation & 50.07 & 63.00 & 50.00 & [36.65, 63.35] & 30 \\
EPIC-KITCHENS-100 & $\geq$1 hand & 90.37 & 92.00 & 89.42 & [83.26, 95.57] & 60 \\
EPIC-KITCHENS-100 & active manipulation & 85.04 & 87.50 & 86.07 & [78.17, 93.97] & 60 \\
\bottomrule
\end{tabular}

\caption{Published figure, judge-only rate, and the gold-corrected PPI++ estimate with its 95\%
interval. Intervals are iid and are lower bounds on their true width, because the evaluation
frames carry no worker identifier to cluster over.}
\label{tab:prev}
\end{table}

Correcting for the error in Table~\ref{tab:errdir} moves every prevalence estimate and, more to
the point, attaches an interval to quantities that are published without one
(Table~\ref{tab:prev}).

The comparative claim is what the vendor actually sells, and it is a difference rather than a
rate. Table~\ref{tab:margins} corrects both sides of it.

\begin{table}[t]
\centering\small
\begin{tabular}{@{}llrrlc@{}}
\toprule
vs. & task & published & corrected & 95\% CI & published inside \\
\midrule
EPIC-KITCHENS-100 & $\geq$1 hand & +6.05 & +3.02 & [-5.23, +11.26] & yes \\
EPIC-KITCHENS-100 & active manipulation & +6.62 & -1.02 & [-11.68, +9.65] & yes \\
Ego4D & $\geq$1 hand & +29.09 & +31.23 & [+18.95, +43.50] & yes \\
Ego4D & active manipulation & +41.59 & +35.06 & [+19.91, +50.20] & yes \\
\bottomrule
\end{tabular}

\caption{Published margin against the gold-corrected margin. Arms are disjoint frame sets and
their variances are summed, which is conservative here: a shared rater or judge offset is
positively correlated across arms and would shrink the variance of the difference.}
\label{tab:margins}
\end{table}

The manipulation margin against EPIC-KITCHENS-100 changes sign under correction and is the only
one of the four that does. At $-1.02$ percentage points it is close enough to zero that the sign
carries little information, and the corrected interval covers the published value. We therefore
cannot refute that margin. We can say that the data does not settle it, and the distinction is
not cosmetic: a wide interval containing the null is an underpowered result, not a negative one,
and treating the two as equivalent is the error this audit exists to object to.

This estimand is not pre-registered and is reported as exploratory. The pre-registered version
asked a narrower question, whether judge \emph{error rate} differs by domain, and it is reported
unchanged in Section~\ref{sec:notheld}.

\section{Clustering and the design effect}
\label{sec:cluster}

The result in this section depends on no human label.

Frames from one camera wearer share scene, lighting, task and equipment, so a corpus statistic
computed over frames is not computed over independent observations. The released evaluation
frames cannot show this, so we went to the raw corpus: 19{,}495 WebDataset shards holding 192{,}903
clips, roughly 16\,TB. A tar header is 512 bytes at a computable offset, so an index of every clip
can be built from HTTP range requests alone, and one frame per clip can then be extracted from an
h265 MP4 inside a tar over the same mechanism. Indexing every shard took 23 minutes and moved
about 150\,MB.

We drew 10{,}000 frames twice, once uniformly over frames and once stratified by factory
worker-hours with at most one frame per clip, and compared a cluster bootstrap over worker
identifier against its iid counterpart.

\begin{table}[t]
\centering\small
\begin{tabular}{lrr}
\toprule
figure & S10k-U (1966 clusters) & S10k-S (1999 clusters) \\
\midrule
$\geq$1 hand & 1.25 & 1.31 \\
2 hands & 1.62 & 1.66 \\
active manipulation & 1.27 & 1.29 \\
\bottomrule
\end{tabular}

\caption{Design effect, defined as the squared ratio of clustered to iid interval width,
$B = 10{,}000$, seed 777.}
\label{tab:deff}
\end{table}

We pre-registered a design effect of at least 2 and Table~\ref{tab:deff} does not reach it, so
that hypothesis fails. What the measurement does show is that the effect is real and smaller than
we predicted: every one of the six figures exceeds 1, and an interval computed as if these frames
were independent is 12 to 29\% too narrow.

The licensed conclusion is narrow and we do not widen it. The corpus has this property. A
frame-level interval drawn from it inherits some of it. Nothing here measures the vendor's own
draw, and nothing can, because the release ships no grouping variable.

\section{Results that did not hold}
\label{sec:notheld}

\paragraph{Prompt variants agree.} The vendor's dataset card and its shipped prompt file differ in
four places, including a substantive rule about whose hands to ignore. We expected the two to
produce different rates. They differ by 0.32 percentage points.

\paragraph{Prompt sensitivity is real and small.} Across eight variants on 2{,}000 frames,
manipulation rates spread 1.25 percentage points and hand-count rates 0.25, against a
pre-registered bar of 5. The single strongest test, stating explicitly that a gloved hand counts,
moved the hand figure by 0.05 points in a corpus of factory work where gloves are near-universal.

\paragraph{Domain bias points the predicted way and misses the bar.} Judge error on manipulation
is 4.76\% on Egocentric frames against 8.33\% on EPIC-KITCHENS, a gap of 3.57 points with EPIC
higher, which is the direction we predicted and below the 5-point threshold we set. At
\vernierNPrimary{} labels it points this way; at 93 it pointed the other way. A hypothesis whose
sign reverses when the sample grows is reporting on the sample. Our own power analysis, run
before any labelling, called this arm underpowered at the pre-registered size, and there is a
confound we disclosed before collecting data and never probed: the judge has plausibly seen
EPIC-KITCHENS and Ego4D in pretraining and cannot have seen the gated commercial corpus, which
pushes accuracy in the opposite direction from the prediction.

\paragraph{Distillation misses its target.} A linear probe on frozen DINOv2 features reaches
0.6933 agreement with the teacher against a 0.90 target, and the abstention cascade cannot hold a
0.80 agreement floor at any coverage above zero at 95\% confidence on a 46-frame calibration
split. The probe is published as a negative result.

\paragraph{Calibration is degenerate by construction.} Expected calibration error is 0.1043 on
hand count and 0.0782 on manipulation, but almost every frame lands in one confidence bin under
greedy decoding. The number describes the decoding strategy more than the judge.

\section{What moved while we watched}
\label{sec:moved}

An audit that publishes only its findings misleads by omission, so this section reports the
corrections we made to our own claims. All five were found by a person reading the prose against
the data, none by an automated check, and the automated check that now exists was written in
response to the first two.

\paragraph{The blind re-label did not wait.} Our protocol specifies a re-label at least seven days
after the first pass, so that the check measures whether the rubric is decidable rather than
whether the rater remembers the frame. Measured from the labels' own timestamps, the median
separation is 2.4 hours and no pair meets the rule. Our writeup had asserted the seven days as
fact. The agreement figures do not change; what they license does, and they are now reported as
within-session consistency.

\paragraph{A control arm arrived after the labels it controlled.} We re-reviewed the frames where
judge and rater disagreed, mixing in an equal number of frames where they agreed so that revision
rates could be compared. The control arm was added in a later commit than the labels, which had
been collected against a plan containing disagreement frames only, so the pass ran unblinded.

\paragraph{A rubric rule nothing enforced.} Rule 12 states that zero visible hands implies no
manipulation. Seven collected labels violated it and neither the data model nor either labelling
tool rejected the combination. Correcting them moved judge-human agreement, moved two prevalence
estimates, and removed a clean zero-errors result on one arm that an earlier draft had been built
on.

\paragraph{A coverage table claimed an experiment we never ran.} The document listing what the
audit does \emph{not} cover reported the cross-domain accuracy experiment as tested. That
experiment specifies a model with cluster-robust intervals. It was never fitted; what exists is a
comparison of two raw error rates.

\paragraph{Buying labels made the vendor's number look better.} Having priced the unresolved
margin at roughly 42 gold frames per arm, we collected 60 more. The corrected margins moved
toward the published figures rather than away, the estimate of what it would take to separate
them rose to 59 per arm, and the asymmetric-error result weakened from absolute to 26 of 28. Two
properties of that second batch belong in the record: it was collected after the first results
were public, so the rater knew which direction would move the margin, and it was labelled at a
median of three seconds a frame against ten for the first batch, with 24 of 60 at two seconds or
less.

\section{Limitations}

One rater, \vernierNPrimary{} frames, and no inter-rater agreement statistic anywhere in this
work, because there is one person. Every result that depends on human gold inherits that. The
blind re-label is a weaker substitute and, as executed, weaker again for the reason given in
Section~\ref{sec:moved}. The rater wrote the rubric after reading the vendor's prompt, which is
anchoring and is not removable at $n=1$. Sixty of the labels were collected after results were
public and at a third of the earlier pace.

The two judge arms are not independent opinions: the open judge and the vendor's stored labels
plausibly share pretraining data, so their agreement is an upper bound on reliability rather than
a measure of it. The design effect is measured on our draws and never on the vendor's evaluation
frames, and cannot be until a grouping variable ships. Every interval on the evaluation arms is
iid and is a lower bound on its true width. The margin estimand is exploratory. Half of the
advertised guarantee, a pose-accuracy claim, requires annotations no public release ships and is
untouched here. Whether hand visibility predicts training value, which is the premise of the
whole metric, is untested.

\section{What a producer could publish}

Four things, each cheap relative to building the corpus. An interval on the headline figure. A
gold sample, with the \emph{direction} of judge error reported rather than only its magnitude,
since a one-directional error makes the published figure an upper bound. A stable frame
identifier that survives between releases, which costs nothing and is what makes an independent
check possible at all. And the exact prompt, versioned, since the card and the shipped file
already differ.

The instrument that produces the range in this paper costs about \$9 per 10{,}000 frames on a
self-hosted open-weights model and needs no API vendor.

\section{Ethics and reproducibility}

This work republishes no frame. The audit's public artifact ships identifiers, labels and
statistics; a small set of downscaled stills appears only in an interactive viewer, chosen so
that a judge-versus-rater disagreement can be inspected, and only from the corpus whose owner
licensed redistribution. Frames showing any person other than the camera wearer are withheld
after a manual review, and that review is by eye and cannot be automated. Adding labels during
this work created 30 new candidate frames that no one had reviewed; they are excluded and the
build refuses to include them.

Worker identifiers are used only as statistical clusters and no per-worker result is reported. We
make no claim about the consent instrument behind the corpus, which is not published, and we
draw no inference about intent from the vendor's repository history; where we report a schema
change between releases we report the fact and stop.

The protocol, the decision log, every result file and the labels are public. The measurement card
reports \texttt{\vernierVerdict} with \vernierNClaims{} claims because one pre-registered item
could not be checked at all. This paper was built from commit \texttt{\vernierGitRev} against
measurement card digest \texttt{\vernierCardDigest}, and every table in it is generated from
those artifacts rather than transcribed.

\section{Conclusion}

A published corpus quality statistic produced by one model pass is an upper bound with an
unstated width. On this corpus the width is at least 12 to 29\% wider than an independent-frames
calculation would suggest, and the direction of the bias is knowable and consistent across two
model families. The comparative claim that the corpus leads a well-known academic dataset does
not survive correction cleanly and does not fail either, and settling it needs about 59 human
labels per arm rather than a better model.

\bibliographystyle{plainnat}
\bibliography{refs}

\end{document}
